\documentclass[10pt]{article}

\usepackage[T1]{fontenc} % Allows pipe character and other fonts
\usepackage[dvipsnames]{xcolor} % Allow Colour
\usepackage[hidelinks]{hyperref} % Allows links
\usepackage[margin=0.2in]{geometry}
\usepackage{enumitem} % Change enumeration spacing etc...
\usepackage{multicol} % Allow multicolumns
\usepackage{titlesec}
\usepackage{titling}
\usepackage{tgbonum}


\titleformat{\section}
{\large\bfseries\color{black}} % Section colour
{}
{0em}
{}[\color{black}\titlerule]

\titleformat{\subsection}
[runin] % Text on same line
{\bfseries}
{}
{0em}
{}

\titleformat{\subsubsection}
{\bfseries}
{}
{1em}
{}

\titlespacing\section{0pt}{8pt plus 4pt minus 2pt}{4pt plus 2pt minus 2pt}

\titlespacing\subsection{0pt}{4pt plus 2pt minus 2pt}{4pt plus 2pt minus 2pt}

\titlespacing{\subsubsection}{0em}{1em}{0em}

\setlist[itemize]{noitemsep, topsep=0pt} % Globally remove spacing between itemize
\raggedcolumns % Turn off stretching columns to fill after columnbreak

\renewcommand{\familydefault}{\sfdefault} % Default sans serif font
\renewcommand{\maketitle}{
    \begin{flushleft}{\Huge\bfseries\theauthor}
    \end{flushleft}

    \begin{center}
        \vspace{-8pt}
        michaelkwan1238@gmail.com | \href{https://github.com/Phys09}{github.com/Phys09} | 647 882 8209 | Richmond Hill ON L4E3X5
        \vspace{2pt}
        \rule[9pt]{\textwidth}{1.5pt}
        \vspace{1pt}
    \end{center}
}

\renewcommand\labelitemi{-} % Rename what is put in front of itemize

\pagenumbering{gobble} % Disable Page Numbering


% ====== BEGIN DOCUMENT ======
\begin{document}

\title{R\'esum\'e}
\author{
  \fontsize{16}{4pt}\selectfont Michael Kwan
}

\maketitle
\vspace{-40pt}

\section{Skills}
\subsection{Languages:}
Python | Java | Javascript/Typescript | C | C\# | HTML | LaTeX | Markdown | R
\subsection{Libraries/Frameworks:}
Flask | Sqlite3 | React | JavaFX | Numpy
\subsection{Tools:}
Git | Visual Studio | Vim | Eclipse | Google Collab | Jupyter | RStudio | Pycharm
\subsection{Operating Systems:}
Windows | Linux/Unix (Ubuntu)

\subsection{Familiarity With:}
Neo4j | PostgreSQL

\section{Education}

\subsection{University of Toronto Mississauga, Honours Bachelor of Science - Mathematics (2018 - Present)}
\begin{itemize}
  \item Relevant courses in Software Engineering, Software Design, Data Structures and Algorithms, Systems Programming, \\
        Theory of Computation, Computer Organization.
        % TA APPLICATION: \item Computer Science courses: CSC108, CSC148, CSC207, CSC209, CSC263, CSC301, CSC302, CSC311 % Remember to Remove!
\end{itemize}

% ===== BEGIN TA APPLICATION =====
%
% \section{Teaching Related Experience:}
% \subsection{UTM Simpledoc Project}
% \begin{itemize}
%     \item Contributed towards the University of Toronto Mississauga Simpledoc project, which had the goal of simply explaining introductory python + computer science concepts to first year students.
%     \item Created a page on trees which introduced its definition, the types of traversal algorithms, and how they can be implemented in python.
%     \item Showed ability to introduce popular data structures such as trees, through the creation of the page.
% \end{itemize}
%
% \subsection{Introduction to Technology}
% \begin{itemize}
%     \item Introduced various concepts related to technology such as how programs work, how they run, and how to configure them to an elderly group of people.
%     \item Guided individuals 1-on-1 on how to setup their computers, allowing them to run programs of their choosing and remember how it was done.
%     \item Increased independence of those who were taught by reducing their reliance on others to help them with technology-related issues.
% \end{itemize}
%
% \section{Course Achievements}
% \subsection{Introduction to Computer Science CSC148H5}
% \begin{itemize}
%     \item Demonstrated proficiency in python and object oriented programming with an understanding of core data structures
%           and complexity, achieving a grade of 81\% and 80\% in the respective courses.
%     \item Developed a python top-down RPG game using the library pygame using the recently taught object oriented concepts that was well documented and tested.
% \end{itemize}
%
% ====== END OF TA APPLICATION =====
\section{Team-Related Development Experience}

%     \subsection{House Price Estimator (Python, Numpy)}
%     \begin{itemize}
%         \item Worked on a program that would take existing training data, to then estimate future house prices
%         using k-nearest-neighbours, and linear regression.
%     \end{itemize}

% Sigmatech banking application with Abtin, Peter + the boys
\subsection{SigmaTech Banking Application (Javascript, CSS, React, PostgreSQL) \hfill Winter 2022}
\begin{itemize}
  \item Collaborated in a team of 7 to create the frontend and backend of a banking web application following the SCRUM software development life cycle.
  \item Built with a React frontend, and a PostgreSQL + Nodejs backend, allows users to manage account information, transfer money, apply for loans, and more.
        Contributions include the creation + design of webpages to interact with database endpoints designed by the backend team.
        % -- Unnecessary fluff -- \item Made to further deepen knowledge and understanding of how react web applications are made, and how databases can be integrated seamlessly into the web app.
\end{itemize}

% CSC302 Final Assignment with group members jeremiah, and menghao. Data visualization tool
\subsection{Data Visualization Tool (Python, Flask, React, CSS) \hfill Fall 2022}
\begin{itemize}
  \item Collaborated in a 3-person team to create a data visualization react web app that displays a dataset on a graph.
  \item With a backend server written in python using flask and sqlite3, data was taken from a dataset, and added to a SQL database to be later accessed through endpoints.
  \item Contributions consist of the creation of endpoints on the backend flask server allowing the frontend to display the data stored in the database.
\end{itemize}

\subsection{Careconnect (Javascript, React, CSS) \hfill Spring 2023}
\begin{itemize}
  \item Built with React and Javascript in a team of 4 for the 2023 Deerhacks hackathon, created a website which pairs users with a personal support worker (PSW) based on needs, location, language, and situation.
  \item Responsibilities included the creation + formatting of pages, algorithm debugging + testing.
\end{itemize}

\subsection{Othello Board Game (Java, JavaFX) \hfill Fall 2019}
\begin{itemize}
  \item With Java, implemented the game Othello in a team of 5 following the SCRUM software development cycle.
  \item Produced a full GUI using the JavaFX library, filled with animations where players could play a game vs other players, or choose from multiple ai opponents.
\end{itemize}

\section{Other Experience}
% CSC301 Assignment 1, using dagger for dependency injection to implement the 6 degrees of Kevin bacon
\subsection{6 Degrees of Kevin Bacon (Java, Neo4j, Dagger2) \hfill Winter 2022}
\begin{itemize}
  \item Implemented the 6 degrees of Kevin Bacon, which connected the actor Kevin Bacon to other actors through at most 6 actors.
  \item Written in Java, the program used dagger2, a dependency injection framework to supply the components of the program. These included the neo4j
        data-access object, the server, and the request handlers.
\end{itemize}

\subsection{Simulated File System (C) \hfill Winter 2020}
\begin{itemize}
  \item Using the C programming language, and binary I/O functions, created a simulated filesystem where a user would be able to create, delete, and search for files in the binary file database.
  \item The program stored information in a binary file, and the program would modify, or read the file to perform operations, similarly to an actual file system.
  \item Gained a better understanding of how computers read/write data to and from files, how Linux file systems store data, and how they are accessed.
\end{itemize}

\subsection{Paint Application (Java, JavaFX) \hfill Fall 2019}
\begin{itemize}
  \item Using object-oriented design patterns (Model-view-controller, Factory) created a java paint application where users could paint on the screen using different brushes/tools.
  \item Created a complete GUI using the JavaFX library.
\end{itemize}

\subsection{Pythonic RPG (Python) \hfill Summer 2019}
\begin{itemize}
  \item Using python 3, and the pygame library, created a top-down RPG game where the player controlled a character to complete levels by completing objectives, and getting to the end.
  \item The pygame library was used to create game objects such as walls the player needed to move around, enemies to avoid, objectives to collect, and items to push around.
\end{itemize}
\end{document}
